\section{Course Outline}

Below is a list of topics and the corresponding textbook chapters.

\begin{table}[h]
\centering
\begin{tabular}{cllp{4cm}}
\toprule
Week & Lesson \# & Topic & Reading \\
\midrule
1 \& 2  & \textbf{1} & Review and Skill Building & Ch. 1, 2 \& Appendix 5 \\
3 \& 4  & \textbf{2} & Acid Equilibrium, Kinetics \& Catalysis & Ch. 5, 7 \& 9 \\
5 \& 7  & \textbf{3} & Linear Free-Energy Relationships & Ch. 8 \\
6        &           & Break Week & \\
8        & \textbf{4} & Test \#1 \& Isotope Effects & Ch. 8 \\
9 \& 10 & \textbf{5} & Substitution and Rearrangement Mechanisms & Ch. 11 \& literature \\
          &           & \textbf{All topics below can be changed. What do you want?} & \\
11 \& 12 & \textbf{6} & Addition and Elimination Mechanisms \& Test \#2 & Ch. 10 \& literature \\
13       & \textbf{7} & Pericyclic Reactions & Ch. 15 \& literature \\
14       & \textbf{8} & Review and Wrap-up & literature \\
\bottomrule
\end{tabular}
\end{table}

\bigskip To achieve our goals we must examine the following ideas.

\subsection{The Past is Prologue}

We will begin with some of the \textbf{core ideas of chemistry}, focusing on structure and energy. This will not be merely a review; it will be a \textbf{deep dive}. There are things we did not tell you. You are ready at last.

We must first start with an understanding of molecular structure. We will consider the \textbf{nature of the chemical bond}. It is not a line on a page, it is an arrangement of electrons and nuclei that form a stable system. Changes in structure can affect the energy and therefore the reactivity. \textbf{Structure affects reactivity}.

To fully understand how structure influences reactivity, we must understand the path from reactant to product. We will develop a deeper understanding of the \textbf{reaction coordinate} and the \textbf{transition state}. We will see how changing the structure of reactants can change the timing of bond changes in the transition state. We will express this via two-dimensional reaction coordinate diagrams and three-dimensional reaction \textbf{energy surfaces} as we try to get as close to the truth as simple line drawings will permit.

\subsection{Seeing Farther}

Now we will move into topics that are familiar but include new ideas for interpreting how structure affects reactivity. We will review the mathematics of reaction kinetics and rate laws and apply these to catalyzed reactions.

We will encounter \textbf{acid/base chemistry} yet again. The fifth time is the charm. The mathematics of acid equilibria will be explored and we will leave the world of water behind and experience a place where strong acids show their true power. Can you \textbf{measure} the \emph{pH} of pure sulphuric acid? Physical organic chemistry will show you how.

\subsection{Structure and Reactivity}

The heart of physical organic chemistry is the \textbf{correlation} of one system with another. Changing the structure of substituted benzoic acids changes the \emph{pK\textsubscript{a}} value. Can we measure the \textbf{electronic effect} of these structural changes using \emph{pK\textsubscript{a}} values as the score? Can we score the \textbf{size} of chemical groups by how they affect the rate of a given reaction? Can we then use these scores to study the \textbf{mechanism} of unrelated systems? Yes we can.

Free energy changes in one system may relate to free energy changes in another system. These correlations are called \textbf{linear free energy relationships (LFER)} and represent one of the most important contributions of physical organic chemistry.

If we can quantify how structure affects reactivity, perhaps we can quantify how structure affects the activity of drugs. This is the idea of \textbf{quantitative structure--activity relationships (QSAR)} and is an important tool in drug development.

Even tiny changes can have measurable effects. Exchanging a hydrogen for \textbf{deuterium} may seem insignificant, yet it can change the activation energy of a reaction. Using \textbf{isotope effects}, we can investigate reactions through changes that are neither structural nor electronic. This idea is called the \textbf{kinetic isotope effect}.

We can \textbf{interpret} the details of reaction mechanisms through observation and detective work. Many \textbf{clever methods} have been developed, and we will explore several of them.

\subsection{Molecular Orbitals}

To investigate the effect of structural changes on electronic arrangements in molecules, we must consider the quantum nature of chemical systems.

We will use simple linear combinations of \textbf{atomic orbitals} to build graphical models of \textbf{molecular orbitals} in small molecules. We will then use simple mathematical approaches to estimate \textbf{energies and electron densities} in $\pi$ systems and explore \textbf{electron delocalization and aromaticity}.

\subsection{Explorations \& Examples}

\textbf{We can go anywhere.} What would you like to investigate?

Each year we will create a body of work that future students can build upon. We will begin with \textbf{examples} of physical organic chemistry being used to investigate \textbf{interesting mechanistic ideas}. We may explore addition and elimination mechanisms, substitution and rearrangement mechanisms, and pericyclic reactions, \textbf{unless you have other ideas}.

We will \textbf{go wherever you wish}. Would you like to explore physical organic chemistry as it is applied to a \textbf{particular problem}? Would you like to \textbf{survey the work} of an important physical organic chemist? We can do anything, and you will document it all in the form of a report.

You will work on your personal exploration of a \textbf{topic} that interests you while we \textbf{explore} examples of reaction mechanisms described above.


%\newpage
\section{Lecture Plan}

Below is a proposed plan for the course this year. \textbf{This will change} as I get to know the preferences and interests of the students and as class cancellations occur.


\begin{longtable}{cp{0.8cm}p{1.8cm}p{7cm}p{3cm}}
\toprule
Lesson & Class & Date & Topic & Reading \\
\midrule
\textbf{1} & 1 & Sept. 9 & Welcome and Course Introduction, Review of Structure and Sterics & Ch. 1, 2.1, 2.3 \\
 & 2 & Sept. 11 & Review of Structure and Molecular Orbitals & Ch. 1.2, 1.3, 2.4 \\
 & 3 & Sept. 14 & Review of Mechanisms and Arrow Pushing & Appendix 5 \\
 & 4 & Sept. 16 & Review of Transition State Theory and Reaction Coordinates & Ch. 7.1--7.3 \\
 & 5 & Sept. 18 & Review of Energy Surfaces & Ch. 7.8, Appendix 5 \\
\textbf{2} & 6 & Sept. 21 & Brønsted Theory and Acidity Functions & Ch. 5.1--5.3 \\
 & 7 & Sept. 23 & Structure Effects on Acid Equilibrium & Ch. 5.4 \\
 & 8 & Sept. 25 & Reaction Kinetics and Acid/Base Catalysis & Ch. 7.4--7.6, 9.1, 9.2 \\
 & 9 & Oct. 2 & Brønsted Plots & Ch. 8.5, 9.3 \\
 & 10 & Sept. 28 & pH-Rate Profiles & Ch. 9.2, 9.3 \\
\textbf{3} & 11 & Oct. 5 & LFER: Inductive Substituent Effects & Ch. 8.2, 8.3 \\
 & 12 & Oct. 7 & Resonance Substituent Effects & Ch. 8.2, 8.3 \\
 & 13 & Oct. 9 & Interpreting Hammett Plots & Ch. 8.2, 8.3 \\
 & & & \textbf{Break Week} & \\
 & 14 & Oct. 19 & Steric Substituent Effects & Ch. 8.4 \\
 & 15 & Oct. 21 & Solvent Effects & Ch. 8.4--8.6 \\
 & 16 & Oct. 23 & Investigating Reaction Mechanisms & Ch. 8.7, 8.8 \\
 & 17 & Oct. 26 & \textbf{Test \#1} & Classes 1--17 \\
 
 
\textbf{4}&18    &Oct. 28  & Theory of isotope effects                                 &Ch. 8.1                \\
      &19    &Oct. 30  & Interpreting isotope effects                              &Ch. 8.1                \\
\textbf{5}&20    &Nov. 2  & Examples of $S_N1$ \& $S_N2$ substitution  &Ch. 11.1 to 11.5, lit.   \\
      &21    &Nov. 4  & Examples of carbocation rearrangements                     &Ch. 11.8, lit.           \\
      &22    &Nov. 6  & Examples of nucleophilic rearrangements                    &Ch. 11.9, 11.10, lit.    \\
      &23    &Nov. 9   & Example exploration: Non-classical cations                &Ch. 11.5, 14.5, lit.     \\
      &      &Nov. 11   & \textbf{No Class} -- Remembrance Day                             &     \\
      &24    &Nov. 13   & Example exploration: Non-classical cations               &Ch. 11.5, 14.5, lit.  \\
      &      &         & \textbf{All topics below can be changed.}   &              \\
\textbf{6}&25    &Nov. 16  & Examples of pH-rate profile for aspirin                   &Ch. 9.1 to 9.3, 10.17 \\
      &26    &Nov. 18  & Examples of acetal formation/hydrolysis                   &Ch. 10.2, 10.12, lit.     \\
      &27    &Nov. 20  & Examples of electrophilic addition  &Ch. 10.1, 10.3 to 10.5\\
      &28    &Nov. 23  &\textbf{Test \#2}                                              &Classes 18 to 29        \\
      &29    &Nov. 25  & Examples of enzyme catalysis                              &Ch. 9.4, 10.2, 10.12   \\   
      &30    &Nov. 27  & Examples of a rearrangement                               &Ch. 11.8 to 11.10, lit. \\
\textbf{7}&31    &Nov. 30  & Examples of cycloaddition reactions                       &Ch. 15.1 to 15.3, lit.   \\
      &32    &Dec. 2   & Examples of electrocyclic reactions                       &Ch. 15.4, lit.           \\
      &33    &Dec. 4   & Examples of pericyclic rearrangements                     &Ch. 15.5, 15.7, lit.     \\
      &34    &Dec. 7  & Examples of carbene reactions                              &Ch. 15.6, lit.           \\
      &35    &Dec. 9  & Wrapup and reflection                                      &           \\





\bottomrule
\end{longtable}


%\newpage
\section{Assessment Plan}

Below is a plan for the \textbf{assessments} in this course.

\begin{longtable}{cp{2cm}p{2cm}p{7cm}c}
\toprule
Assignment & Due Before & Due Date & Objective & Grade \\
\midrule
0 & Meeting 2 & Sept. 11 & First literature exploration \& Learning Plan & 4\% \\
1 & Meeting 5 & Sept. 18 & Organic chemistry review & 4\% \\
2 & Meeting 8 & Sept. 25 & Dogma of physical organic chemistry & 4\% \\
3 & Meeting 11 & Oct. 5 & Acid/base catalysis and Brønsted plots & 4\% \\
4 & Meeting 13 & Oct. 9 & Hammett plots & 4\% \\
\textbf{Test \#1} & Meeting 17 & Oct. 26 & Covers meetings 1--16 & 15\% \\
5 & Meeting 19 & Oct. 30 & Transition-state structure & 4\% \\
6 & Meeting 24 & Nov. 13 & Data analysis and mechanism interpretation & 4\% \\
7 & Meeting 26 & Nov. 18 & Personal exploration proposal & 4\% \\
\textbf{Test \#2} & Meeting 28 & Nov. 23 & Covers Meetings 18--27 & 15\% \\
8 & Meeting 30 & Nov. 27 & Data analysis and mechanism interpretation & 4\% \\
9 & Meeting 33 & Dec. 4 & First draft of personal exploration & 4\% \\
\textbf{Report} & Meeting 35 & Dec. 9 & Submit final personal exploration report & 20\% \\
\bottomrule
\end{longtable}


